\documentclass[12pt]{article}
\usepackage[utf8]{inputenc}
\usepackage{graphicx}
\usepackage{hyperref}
\usepackage{amssymb}
\usepackage[T1]{fontenc}
\usepackage{lmodern}
\usepackage{textcomp}
\usepackage{pgfplots}
\usepackage{pgfplotstable}
\usepackage{longtable}
\usepackage[capitalise]{cleveref}
\usepackage{caption}
\pgfplotsset{compat=1.18}
\captionsetup[table]{position=top,justification=raggedright,singlelinecheck=false}
\captionsetup[longtable]{position=top,justification=raggedright,singlelinecheck=false}
\captionsetup[figure]{justification=raggedright,singlelinecheck=false}
\setlength\LTleft{0pt}
\setlength\LTright{0pt}
\setlength\LTcapwidth{\textwidth}
\pgfplotstableset{
  show variable once/.style={
    columns/variable/.style={
      string type,
      column name=Variable,
      postproc cell content/.code={
        \pgfkeysgetvalue{/pgfplots/table/@cell content}{\cellcontent}
        \edef\currentvar{\cellcontent}
        \ifx\currentvar\lastvar
          \pgfkeyssetvalue{/pgfplots/table/@cell content}{}
        \else
          \ifnum\pgfplotstablerow>0
            \pgfkeyssetevalue{/pgfplots/table/@cell content}{\noexpand\hline\space\cellcontent}
          \fi
          \xdef\lastvar{\currentvar}
        \fi
      }
    }
  }
}
\def\lastvar{}
\setlength{\parskip}{0.8em}
\setlength{\parindent}{0pt}
%\renewcommand{\thesubsection}{\thesection.\alph{subsection}}

\title{\Large\bfseries ASSIGNMENT 2: REPEATED MEASURES ANOVA AND
MANOVA}
\author{Chenghao Wang}
\date{\today}

\begin{document}

\maketitle

\section{Descriptive Analysis}
\subsection{Descriptive statistics}\label{subsec:descriptive-statistics}
We find many missing values in PostBD FEV1 (the continuous outcome variable) 
in the dataset. Since both repeated measures ANOVA and MANOVA are complete case analyses, 
we filtered out time points with $\ge 20\%$ missing values in PostBD FEV1, and then
we only kept subjects whose PostBD FEV1s have no missing values at the retained time points.
471 out of 695 subjects ( $\sim 68\%$) and
13 out of 20 time points were included in the analysis for this assignment.
Across the treatment group, 67-70\% of subjects were retained.

\pgfplotstabletypeset[
  col sep=comma,
  ignore chars={"},
  begin table=\begin{longtable},
  end table=\end{longtable},
  columns={visitc,mean_budesonide,sd_budesonide,mean_nedocromil,sd_nedocromil,mean_placebo,sd_placebo},
  columns/visitc/.style={column name=\shortstack{Visit\\(month)}, fixed, precision=0},
  columns/mean_budesonide/.style={column name=\shortstack{Mean\\(Bud)}, fixed, fixed zerofill, precision=2},
  columns/sd_budesonide/.style={column name=\shortstack{SD\\(Bud)}, fixed, fixed zerofill, precision=2},
  columns/mean_nedocromil/.style={column name=\shortstack{Mean\\(Ned)}, fixed, fixed zerofill, precision=2},
  columns/sd_nedocromil/.style={column name=\shortstack{SD\\(Ned)}, fixed, fixed zerofill, precision=2},
  columns/mean_placebo/.style={column name=\shortstack{Mean\\(Plbo)}, fixed, fixed zerofill, precision=2},
  columns/sd_placebo/.style={column name=\shortstack{SD\\(Plbo)}, fixed, fixed zerofill, precision=2},
  every head row/.style={
    before row=\caption{PostBD FEV1 mean and SD by treatment group and visit}\label{tab:descriptive-summary}\\\hline,
    after row=\hline\endfirsthead
      \multicolumn{7}{l}{\tablename\ \thetable\ (continued)}\\\hline
      \shortstack{Visit\\(month)} & \shortstack{Mean\\(Bud)} & \shortstack{SD\\(Bud)} & \shortstack{Mean\\(Ned)} & \shortstack{SD\\(Ned)} & \shortstack{Mean\\(Plbo)} & \shortstack{SD\\(Plbo)}\\\hline\endhead
      \hline\endfoot
      \hline\endlastfoot
  }
]{../hw2_code/output/descriptive_summary.csv}

\subsection{Interaction plot}
\begin{figure}[ht]
\centering
\includegraphics[width=0.85\textwidth]{../hw2_code/output/interaction_plot.png}
\caption{Mean PostBD FEV1 over time by treatment group}
\label{fig:interaction-plot}
\end{figure}

Since we performed repeated measures ANOVA and MANOVA, 
we treated time
as a discrete variable.
As shown in \cref{fig:interaction-plot}, the profiles do not appear parrallel,
and the treatment groups appear to differ in their overall level.

\section{Univariate Repeated Measures ANOVA}

All the analysis was performed using R for this assignment.
Nonetheless, we want to highlight the potential discrepancy between
the results from R and SAS and the corresponding reasons, 
as this might help us better
understand the technical details of repeated measures
ANOVA and MANOVA. The behavioral differences between R and SAS 
in this case
were identified using the dental study data from class, 
and the relevant comparisons are not reported here.

By default, for categorical variables, R uses reference coding, whereas SAS uses sum coding.
Choice of coding scheme does not affect fitted values and Type I SS, but it can affect Type III SS.
In this assignment, we set the coding scheme as sum coding, and we use Type III SS, with \texttt{car} package,
except one test of group effect based on MANOVA, where \verb|stats::manova| is used.

\subsection{Sphericity assessment}


The statistics of Mauchly's test for sphericity from R and SAS are identical.
Yet the p-values are different between R and SAS, because
their ways of approximation differ. The adjusted p-values are the same.

For this univariate repeated measures ANOVA, I used PostBD FEV1
as the response variable, with treatment group (between-subject)
and visit time (within-subject) as
explanatory variables.
I used the filtered data from \cref{subsec:descriptive-statistics}.

The Mauchly's test for sphericity shows a significant
violation of the sphericity assumption 
(test statistic = $8 \times 10^{-7}$ ; p < 0.001).
Therefore, we do not have a compound symmetry covariance structure,
and thus we no longer have exact F tests for within-subject effects (visit time,
treatment group $\times$ visit time). In this case, we should
report the adjusted p-values rather than unadjusted p-values
for the within-subject effects.
It is worth mentioning that the Mauchly's test for sphericity
is indeed for evaluating whether the covariance structure is
of Type H, of which compound symmetry is a special case. 


\subsection{Hypothesis tests}
\subsubsection{Group $\times$ Time interaction}
The F statistic for treatment group $\times$ visit time interaction effect
is 1.34.
The Greenhouse-Geisser (G-G) adjusted p-value for the 
interaction effect is 0.245,
and the Huynh-Feldt (H-F) adjusted p-value for the interaction effect is 0.244.
Both indicate that the interaction effect is insignificant.
Thus we conclude that the PostBD FEV1 mean trajectories for the three treatment groups
are parallel.
Incidentally, the unadjusted p-value for the interaction effect is 0.123,
which is also insignificant.
Since the data is from a randomized trial, the test for
parallelism tells us wether there is evidence of a treatment effect,
if filtering out missingness does not introduce any bias.
Although the test result seems to contradict our findings in \cref{fig:interaction-plot},
there are several possible explanations for this.
First, the unparallel pattern might be due to noise. Based on \cref{tab:descriptive-summary},
the standard deviations are larger than the difference in the means 
between the treatment groups.
Besides, the trajectories appear parallel in some time intervals
while unparallel in others. Yet the test for interaction effect
evaluates the group difference across all time points. Despite the 
local non-parallelism, the global result is slightly affected.

\subsubsection{Time effect}
Since the treament groups have parallel profiles, we directly
tested the visit time effect. The F statistic for visit time effect is 2274.41.
Both the G-G adjusted p-value and the H-F adjusted p-value are < 0.001.
This suggests that the mean PostBD FEV1 changes over visit time.
Incidentally, the unadjusted p-value for visit time effect is also < 0.001.

\subsubsection{Group effect}
With parallelism, we tested the treatment group effect
averaged over time using an F test,
which indicates an insignficant result
(F = 0.84, p = 0.434). Thus there is no evidence that the treatment groups
differ in their overall mean PostBD FEV1.

\section{MANOVA}
In MANOVA, there are four types of tests for an effect.
The test statistics are as follows:
Wilks' Lambda, Pillai's Trace, Hotelling-Lawley Trace, and Roy's Greates Root.
Each of the four tests gives the same result in R and SAS, except Hotelling-Lawley Trace,
because the approximation differs between R and SAS.
In this assignment, we only report 
Wilks' Lambda and the corresponding p-value.


\subsection{Test of parallelism}
The Wilks' Lambda for 
treatment group $\times$ visit time effect is 0.89 (F = 2.26),
and the corresponding p-value is < 0.001. Thus, we do not reject the null hypothesis
and conclude that the treatment group profiles are parallel.
This shows that the difference in PostBD FEV1 between the treatment groups
do not change over time. Since the data is from a randomized trial,
this test result indicates that there is no evidence of a treatment effect,
if filtering out missingness does not introduce any bias.

\subsection{Test of group effect}\label{subsec:test-group}
We performed two tests in this subsection.
The first test assesses treatment group effect at each time point
simultaneously, which does not rely on the assumption that
the group profiles are parallel. Unlike other tests in this assignment,
this test is performed using
\verb|stats::manova| in R. The Wilks' Lambda for treatment group effect
is 0.86 (F = 2.72), and the corresponding p-value is < 0.001. Thus,
we reject the null and conclude that the postBD FEV1 means
are different across treatment groups at certain visit times.

The second test assesses treatment group effect averaged over time,
which relies on the parallelism assumption (a reasonable assumption
in our case, given the test results for parallelism).
The Wilks' Lambda is 1.00 (F = 0.84), and the corresponding p-value
is 0.434. 
Thus, we do not reject the null and conclude that the
postBD FEV1 means averaged over time
are the same for each treatment group.


\subsection{Test of time effect}
Like \cref{subsec:test-group},
we performed two tests of time effect.
The first test evaluates whether the trajectories
of postBD FEV1 trajectories
are flat for each treatment group, which does not
rely on the parallelism assumption.
The Wilks' Lambda is 0.06 (F = 61.30), and the corresponding p-value
is < 0.001. Thus, we reject the null and conclude that
postBD FEV1 mean trajectories are not flat for certain treatment groups.

The second test evaluates whether the trajectories
of postBD FEV1 mean are flat by averaging over the
treatment groups. The Wilks' Lambda is 0.06 (F = 552.57), and the
corresponding p-value is < 0.001. Thus, we reject the null and conclude that
postBD FEV1 mean trajectories are not flat by averaging over the treatment
groups.

The test of time effect using univariate repeated measures ANOVA
produces a significant result,
which is the same as the test results using MANOVA.
Although the F statistics and the p-values
are different.


\section{Comparison and Interpretation}

\subsection{Comparison of approaches}
The results from the univariate repeated measures ANOVA and
MANOVA lead to the same conclusions,
except the test simultaneously assessing treatment group effect
at each time point in MANOVA, which gives a signficant result.
The test of treatment group effect in univariate repeated measures ANOVA
and the second test of treatment group effect in MANOVA both give
an insignificant result, since these two tests 
evaluate the treatment group effect averaged over time.
The null hypotheses are closely related but not identical.
Besides, as is illustrated in \cref{fig:interaction-plot},
the group profiles appear coincident at some time points while
not coincident at the others, which likely contributes to
the difference in the treatment group effect test results.

Moreover, although most tests produce results that lead to
the same conclusions, yet the F statistics and the p-values
are not identical (identical for the test of group
effect averaged over time, as group effect is between-subject). 
This is because the covariance structures in the models
differ. In univariate ANOVA's model, the covariance matrix is
assumed to be
compound symmetry, while it is assumed to be unstructured in MANOVA.
Since the Mauchly's test for sphericity strongly indicates that
the covariance is not compound symmetry, MANOVA is more
appropriate for our data.

\subsection{Scientific interpretation}

Based on the test results above, we find that
there is no significant difference between the
intervention treatments and placebo in PostBD FEV1,
a measure of lung function. 
Additionally, postBD FEV1 significantly increases over time,
indicating that the lung functions of children with asthma
continuously improve.

\subsection{Limitations}
First, the missingness mechanism is not clear.
Moreover, missing data occurred not only among
certain participants, but also at multiple time points,
where some time points showed substantial proportions
of missingness (> 80\%).
Given that both repeated measures ANOVA and MANOVA are complete case analyses,
we filtered out certain time points and participants, as described
in \cref{subsec:descriptive-statistics}.
However, this may compromise efficiency and introduce bias.

Second, we did not adjust for any covariates other than treatment group
and visit time for simplicity. This can introduce bias and reduce power.

Lastly, we analyzed solely the primary outcome, PostBD FEV1, and
did not assess other outcomes. Accordingly, this may 
influence our conclusions regarding the treatment effect.

\end{document}
